\section{Билет 11. Сложение взаимно перпендикулярных гармонических колебаний. Вывод уравнения. Исследовать форму траекторий для разности фаз равной 0 ($\pi$). Фигуры Лиссажу.}

\begin{definition}
Рассмотрим сложение двух гармонических колебаний, направленных вдоль взаимно перпендикулярных осей $x$ и $y$, имеющих одинаковую циклическую частоту $\omega$:
\begin{equation}
    x(t) = A \cos(\omega t), \quad y(t) = B \cos(\omega t + \varphi),
\end{equation}
где $\varphi$ --- разность фаз складываемых колебаний. Результирующее движение материальной точки происходит в плоскости $xOy$, а форма её траектории зависит от амплитуд $A, B$ и разности фаз $\varphi$.
\end{definition}

\begin{theorem}[Вывод уравнения траектории]
Чтобы получить уравнение траектории в явном виде, необходимо исключить из системы параметр $t$. Из первого уравнения:
\begin{equation}
    \cos(\omega t) = \frac{x}{A}, \quad \sin(\omega t) = \pm\sqrt{1 - \frac{x^2}{A^2}}.
\end{equation}
Раскроем косинус суммы во втором уравнении и подставим выражения для $\cos(\omega t)$ и $\sin(\omega t)$:
\begin{equation}
    \frac{y}{B} = \cos(\omega t)\cos\varphi - \sin(\omega t)\sin\varphi = \frac{x}{A}\cos\varphi \mp \sqrt{1 - \frac{x^2}{A^2}}\sin\varphi.
\end{equation}
Перенесём первое слагаемое в левую часть и возведём обе части в квадрат:
\begin{equation}
    \left(\frac{y}{B} - \frac{x}{A}\cos\varphi\right)^2 = \left(1 - \frac{x^2}{A^2}\right)\sin^2\varphi.
\end{equation}
Раскрывая скобки и приводя подобные члены, получаем общее уравнение траектории:
\begin{equation}
    \frac{x^2}{A^2} + \frac{y^2}{B^2} - \frac{2xy}{AB}\cos\varphi = \sin^2\varphi. \label{eq:trajectory_general}
\end{equation}
Это уравнение эллипса, оси которого ориентированы относительно координатных осей произвольным образом.
\end{theorem}

\begin{property}[Случай $\varphi = 0$]
Подставляя $\varphi = 0$ в \eqref{eq:trajectory_general} ($\cos 0 = 1, \sin 0 = 0$), получаем:
\begin{equation}
    \frac{x^2}{A^2} + \frac{y^2}{B^2} - \frac{2xy}{AB} = 0 \quad \Rightarrow \quad \left(\frac{x}{A} - \frac{y}{B}\right)^2 = 0.
\end{equation}
Отсюда $y = \frac{B}{A}x$. Точка движется по прямой линии, проходящей через 1-й и 3-й квадранты. Результирующее движение остаётся гармоническим колебанием вдоль этой прямой с амплитудой $\sqrt{A^2 + B^2}$.
\end{property}

\begin{property}[Случай $\varphi = \pm \pi$]
Подставляя $\varphi = \pm \pi$ ($\cos(\pm\pi) = -1, \sin(\pm\pi) = 0$), получаем:
\begin{equation}
    \frac{x^2}{A^2} + \frac{y^2}{B^2} + \frac{2xy}{AB} = 0 \quad \Rightarrow \quad \left(\frac{x}{A} + \frac{y}{B}\right)^2 = 0.
\end{equation}
Отсюда $y = -\frac{B}{A}x$. Траектория также представляет собой прямую линию, но проходящую через 2-й и 4-й квадранты. Амплитуда результирующих колебаний остаётся равной $\sqrt{A^2 + B^2}$.
\end{property}

\begin{remark}
\textbf{Фигуры Лиссажу.} Если частоты складываемых колебаний не равны, но кратны друг другу ($\omega_x : \omega_y = m : n$, где $m, n$ --- целые числа), траектория замыкается, образуя сложные устойчивые кривые, называемые \textbf{фигурами Лиссажу}. Отношение частот равно отношению числа касаний или пересечений фигуры с прямыми, параллельными осям координат: $\omega_x / \omega_y = n_y / n_x$. Фигуры Лиссажу широко применяются для определения неизвестной частоты по известному эталону с помощью осциллографа.
\end{remark}

\vspace{1cm}

\begin{figure}[htbp]
    \centering
    \begin{tikzpicture}[
        scale=0.8,
        every node/.style={font=\small},
        declare function={
            A = 2; B = 1.5; % Амплитуды (A > B для наглядности эллипса)
        }
    ]

    % --- ВАРИАНТ 1: phi = 0 (Свойство 11.1) ---
    \begin{scope}[xshift=-8cm]
        \draw[->] (-2.5, 0) -- (2.5, 0) node[right] {$x$};
        \draw[->] (0, -2) -- (0, 2) node[above] {$y$};

        % Траектория: Прямая y = (B/A)x
        \draw[blue, thick] (-A, -B) -- (A, B);

        % Точки экстремума
        \fill[red] (A, B) circle (2pt);
        \fill[red] (-A, -B) circle (2pt);

        % Подписи
        \node at (0, -2.5) {$\varphi = 0$};
        \node[align=center, font=\tiny] at (2, 2) {Прямая\\в 1 и 3 кв.};
    \end{scope}

    % --- ВАРИАНТ 2: phi = pi/2 (Эллипс) ---
    \begin{scope}[xshift=0cm]
        \draw[->] (-2.5, 0) -- (2.5, 0) node[right] {$x$};
        \draw[->] (0, -2) -- (0, 2) node[above] {$y$};

        % Траектория: Эллипс (A=2, B=1.5)
        \draw[blue, thick] (0,0) ellipse (2 and 1.5);

        % Оси эллипса
        \draw[dashed, gray] (-2, 0) -- (2, 0);
        \draw[dashed, gray] (0, -1.5) -- (0, 1.5);

        % Подписи
        \node at (0, -2.5) {$\varphi = \frac{\pi}{2}$};
        \node[align=center, font=\tiny] at (2, 2) {Эллипс\\(общий случай)};
    \end{scope}

    % --- ВАРИАНТ 3: phi = pi (Свойство 11.2) ---
    \begin{scope}[xshift=8cm]
        \draw[->] (-2.5, 0) -- (2.5, 0) node[right] {$x$};
        \draw[->] (0, -2) -- (0, 2) node[above] {$y$};

        % Траектория: Прямая y = -(B/A)x
        \draw[blue, thick] (-A, B) -- (A, -B);

        % Точки экстремума
        \fill[red] (-A, B) circle (2pt);
        \fill[red] (A, -B) circle (2pt);

        % Подписи
        \node at (0, -2.5) {$\varphi = \pi$};
        \node[align=center, font=\tiny] at (2, 2) {Прямая\\в 2 и 4 кв.};
    \end{scope}

\end{tikzpicture}
\caption{Траектории результирующего движения при сложении перпендикулярных колебаний с разными разностями фаз $\varphi$.}
\label{fig:lissajous_phases}
\end{figure}

\begin{figure}[htbp]
    \centering
    \begin{tikzpicture}[scale=1.2, every node/.style={font=\small}]

        % --- 1:1 (Окружность/Эллипс) ---
        \begin{scope}[shift={(-3, 3)}]
            % Оси
            \draw[->] (-1.5, 0) -- (1.5, 0) node[right] {$x$};
            \draw[->] (0, -1.5) -- (0, 1.5) node[above] {$y$};
            % Фигура: x=sin(t), y=cos(t) -> Окружность (сдвиг фаз pi/2)
            \draw[blue, thick, domain=0:360, samples=100] plot ({sin(\x)}, {cos(\x)});
            \node at (0, -1.8) {$\omega_x : \omega_y = 1:1$};
            \node[font=\tiny, gray] at (0, -2.1) {(окружность)};
        \end{scope}

        % --- 1:2 (Восьмерка горизонтальная) ---
        \begin{scope}[shift={(3, 3)}]
            \draw[->] (-1.5, 0) -- (1.5, 0) node[right] {$x$};
            \draw[->] (0, -1.5) -- (0, 1.5) node[above] {$y$};
            % Фигура: x=sin(t), y=sin(2t)
            \draw[blue, thick, domain=0:360, samples=200] plot ({sin(\x)}, {sin(2*\x)});
            \node at (0, -1.8) {$\omega_x : \omega_y = 1:2$};
        \end{scope}

        % --- 2:1 (Восьмерка вертикальная) ---
        \begin{scope}[shift={(-3, -2)}]
            \draw[->] (-1.5, 0) -- (1.5, 0) node[right] {$x$};
            \draw[->] (0, -1.5) -- (0, 1.5) node[above] {$y$};
            % Фигура: x=sin(2t), y=sin(t)
            \draw[blue, thick, domain=0:360, samples=200] plot ({sin(2*\x)}, {sin(\x)});
            \node at (0, -1.8) {$\omega_x : \omega_y = 2:1$};
        \end{scope}

        % --- 3:2 (Сложная фигура) ---
        \begin{scope}[shift={(3, -2)}]
            \draw[->] (-1.5, 0) -- (1.5, 0) node[right] {$x$};
            \draw[->] (0, -1.5) -- (0, 1.5) node[above] {$y$};
            % Фигура: x=sin(3t), y=sin(2t)
            \draw[blue, thick, domain=0:360, samples=200] plot ({sin(3*\x)}, {sin(2*\x)});
            \node at (0, -1.8) {$\omega_x : \omega_y = 3:2$};
        \end{scope}

    \end{tikzpicture}
	\caption{Типовые фигуры Лиссажу при разности фаз $\varphi = \pi/2$. Форма траектории определяется отношением частот складываемых колебаний. \textbf{ВАЖНО:} фигуры Лиссажу находятся как бы в квадрате и всегда касаются границ}
    \label{fig:lissajous_figures}
\end{figure}
